\section{The fiscal condition}
\label{sec:fiscal}

This module supports the central claim on its revenue side. If the public sector holds the wage
side of the exposure, the question is whether the tax that reaches the capital replacing that
labor could replace the revenue lost. We build both sides from components and report the answer
under every reading we can defend.

\subsection{What replacement requires}

Not every dollar of displaced wages leaves the wage tax base, because some of the displaced are
re-employed, at a wage on average lower than the one they lost.

\begin{definition}[Retained wage share and the fiscal condition]
\label{def:condition}
Let $R \in [0,1]$ be the share of a displaced wage bill that returns as wages, $\tau_\ell$ the
effective tax rate on wages and $\tau_k$ the effective marginal rate on the capital income that
replaces them. Replacing the lost revenue requires
\begin{equation}
\tau_k \;\ge\; \tau_\ell\,(1 - R) \;\equiv\; \tau_k^{\ast},
\label{eq:condition}
\end{equation}
where $\tau_k^{\ast}$ is the required rate.
\end{definition}

Measured on fourteen biennial vintages of displaced worker data and solved as a fixed point
against the stock-flow identity of Section~\ref{sec:data}, $R = \result{RetainedWageShare}$. Two
defensible readings of $\tau_\ell$ exist, a narrower one from the automation tax literature and a
broader bottom-up one built from federal taxes over wages and salaries, giving
$\tau_k^{\ast}$ between \result{RequiredEasy} and \result{RequiredHard} percent. We carry both and
call the lower one the easier reading. The retained wage share and both bounds were reproduced by
the independent rebuild.

\subsection{Building the rate on the income that shifts}

\begin{definition}[The assembled effective rate]
\label{def:tauk}
Let $\sigma$ be the share of the return that is economic rent, $s$ the share of rents booked in
lower-taxed jurisdictions, $\tau_c$ the combined entity-level rate, $\theta$ the share of
corporate equity held in taxable accounts, $\delta$ the present-value factor for deferral between
accrual and realization, $\tau_s$ the statutory shareholder rate, $b$ the debt share of financing
and $\tau_b$ the effective rate on interest income. With the shareholder layer
$\tau_{sh} = \theta\,\delta\,\tau_s$ and the entity layer $e$ formed from $\tau_c$ and $s$,
\begin{equation}
\tau_k \;=\; \sigma\,\bigl[e + (1-e)\,\tau_{sh}\bigr]
\;+\;(1-\sigma)\,\bigl[(1-b)\,\tau_{sh} + b\,(\tau_b - \tau_c)\bigr].
\label{eq:tauk}
\end{equation}
The first bracket is the rate on rents; the second is the rate on the normal return.
\end{definition}

Two features of Equation~\eqref{eq:tauk} do most of the work and neither is our assumption.
Under full expensing the normal return is relieved at the entity level, so at the margin it bears
only the owner's tax \citep{AcemogluManeraRestrepo2020}, which is why the debt-financed term is
negative throughout the sourced range and why the budget office's measured rate on debt-financed
corporate investment, minus \result{CBODebtFinancedETR} percent, sits inside the interval our
components imply \citep{CBO2014}. The rent share is the second: \citet{Barkai2020} puts it at
\result{RentBarkai} and \citet{KarabarbounisNeiman2014} at zero, two readings of what has happened
to the corporate sector rather than two estimates of one quantity, so we report both and never
average them.

\paragraph{The base, and the case against our own choice.} We take the rate on the marginal flow
of US AI surplus, whoever holds the claim, because the question concerns the next dollar that
moves. Against that, an economy-wide average rate on capital income, around twenty percent for
advanced economies, is measured on a stock mostly neither fully expensed nor easily shifted; if AI
capital comes to look like the rest of the capital stock, that is the relevant rate. Which base
answers which question is a judgment, not a fact.

\subsection{Which government, and the answer on each}

The two sides of Equation~\eqref{eq:condition} have to be measured on the same government. The
labor tax readings are federal: both are built from federal taxes on wages, one from the
automation tax literature and one bottom-up from federal receipts over wages and salaries. The
assembled rate, however, adds an effective state corporate tax to the federal entity rate, so a
comparison between them as they stand mixes jurisdictions. We therefore report two consistent
versions and make the federal one the headline, because the claim this paper makes is about the
federal balance sheet.

On the \emph{federal-only} version the state corporate term is set to zero, giving an assembled
rate of \result{TauKFederalOnly} percent against a required \result{RequiredFedEasy} to
\result{RequiredFedHard} percent. On the \emph{all-government} version the state corporate term
stays and the labor side gains the state and local wage-linked tax, worth
\result{StateLocalLaborAddOn} percent of the wage bill, giving \result{TauKAllGovernment} percent
against \result{RequiredAllGovEasy} to \result{RequiredAllGovHard} percent. Making the
comparison consistent widens the gap on either version, because the state adds more to the labor
side than to the capital side. Under the second rent reading the assembled rate is
\result{TauKKN} percent on either jurisdiction.

This is a scenario result in the strict sense, since two components are swept rather than
sourced; Figure~\ref{fig:tauk} shows that sweep. The construction was reproduced in full by the
independent rebuild.

How far short depends on the base for $\theta$, and the two answers differ enough that we report
both. On our base, $\theta$ is measured over all US equity outstanding, foreign holders included,
because a foreign holder bears close to no US tax at the owner level and that dollar is still in
the surplus. On the alternative base, the same numerator over domestically held stock alone,
$\theta$ rises to \result{CRSTheta} percent and the central rate to \result{TauKCRSCentral}
percent.

The precise statement is the share of a stated grid in which the condition closes. That grid is
the set of swept components, each drawn uniformly over the range given for it in
Section~\ref{sec:data} and the replication package; the share is a property of that grid and of
the uniform weighting, and it is not a probability that the condition closes. On our base it is
\result{PassShareOursEasy} percent against the easier labor tax reading and
\result{PassShareOursHard} percent against the harder; on the domestic-holder base it is
\result{PassShareCRSEasy} and \result{PassShareCRSHard} percent. The corners of the space reach further than those shares
suggest: the analytic maximum is \result{AnalyticMaxOurs} percent on our base and
\result{AnalyticMaxCRS} percent on the other, so both exceed the easier required rate of
\result{RequiredEasy} percent, and the second exceeds the harder \result{RequiredHard} percent as
well. What holds on every base and under every reading is that the central case falls short, and
that under the second rent reading the condition closes nowhere in the swept space.

\begin{figure}[htbp]
\centering
\includegraphics[width=0.84\textwidth]{fig3_capital_tax.pdf}
\caption{The assembled effective marginal rate on income that moves from wages to AI capital,
against the rate required to replace the wage taxes lost. Bars are central values, whiskers the
fifth to ninety-fifth percentile of the swept components; the shaded band is the required rate
across the two readings of the labor tax rate. The two left bars use the taxable-holder share
measured over all equity outstanding, the two right bars the same numerator over domestically held
stock. Scenario, on a construction reproduced by independent rebuild.}
\label{fig:tauk}
\end{figure}

\subsection{What would have to be true for the conclusion to reverse}

A shortfall is only interesting if it survives the assumptions that would close it, so we apply
each of those in turn and then together, all on the all-government version for comparability.
Treating the foreign-held share of equity as bearing a fifteen percent treaty withholding rate at
the owner level, rather than nothing, raises the assembled rate to \result{TauKWithholding}
percent; that treaty rate is a stated assumption and not a measurement. Taking the debt-financed
share from the filers' own books rather than the economy-wide stock ratio raises it to
\result{TauKLowDebtShare} percent, because the debt-financed term is negative. Applying both
together gives \result{TauKJointWeakening} percent, still below the required
\result{RequiredAllGovEasy} percent. The second rent reading moves in the opposite direction, to
\result{TauKKN} percent.

One further assumption works in the paper's favor and belongs here for that reason. The labor
tax rate is economy-wide, while displacement is not uniform across the wage distribution. Using
the payroll and income tax rates by wage quintile already in the release, the effective rate runs
from \result{LaborTaxQuintileLow} percent at the bottom to \result{LaborTaxQuintileHigh} percent
at the top, so displacement concentrated in the bottom quintile would require only
\result{RequiredBottomQuintile} percent rather than \result{RequiredTopQuintile} percent at the
top. Even that lower requirement is above the assembled rate on either jurisdiction.

\subsection{Pass-through, and what it can do to the comparison}

The condition in Definition~\ref{def:condition} carries a coefficient that the headline comparison
sets to one. Written as $\tau_k \ge \tau_\ell (1 - R)$, it assumes that every dollar of displaced
wages reappears as a dollar of taxable US capital income. Making that explicit,
\begin{equation}
\tau_k \, g \;\ge\; \tau_\ell\,(1 - R), \qquad
\tau_k^{\ast} = \frac{\tau_\ell\,(1 - R)}{g},
\label{eq:passthrough}
\end{equation}
with $g$ the taxable US capital income generated per dollar of displaced wages, bounded in
$[0,1]$ with output preserved. Everything reported above is the case $g = 1$, and it stays the
headline because $g$ cannot be sourced. Incomplete pass-through can only widen the shortfall,
never close it: any $g$ below one raises the required rate, and the decomposition of a displaced
wage dollar in Appendix~\ref{app:detail}, which puts $g$ at \result{GLow} to \result{GHigh}, is
scenario arithmetic reported there for that reason.

\subsection{The boundary, and what would have to happen for the condition to close}

The bound $g \le 1$ belongs to the output-preserved case. If AI raises output, a displaced wage
dollar can be associated with more than a dollar of taxable capital income, and $g$ can exceed
one. Stating the condition that way turns a comparison of point estimates into a boundary, which
is the more useful object: Figure~\ref{fig:boundary} draws the fiscal balance
$\tau_k \, g - \tau_\ell (1 - R)$ over the retained wage share and $g$, with the zero contour for
each jurisdiction and each labor tax reading.

Two things can be read straight off it. At the measured retained wage share, the condition closes
only if $g$ reaches \result{GStarFedEasy} to \result{GStarFedHard} on the federal-only version,
or \result{GStarAllEasy} to \result{GStarAllHard} on the all-government one. Every one of those
is above the output-preserved ceiling, so with output preserved the condition does not close at
the assembled rate anywhere on the grid. And the conclusion reverses only in the upper region:
either output rises enough to carry $g$ well past one, or the share of displaced wages that comes
back as wages rises far above what the displaced worker data show.

\begin{figure}[htbp]
\centering
\includegraphics[width=0.84\textwidth]{fig6_boundary.pdf}
\caption{The fiscal balance per displaced wage dollar, $\tau_k \, g - \tau_\ell (1 - R)$, over the
retained wage share $R$ and the taxable capital income generated per displaced wage dollar $g$.
Each line is the zero contour for one jurisdiction and one labor tax reading, at the assembled
capital rate for that jurisdiction; the condition closes above a line and fails below it. The
shaded region is where it fails on the headline federal-only version. The horizontal line is the
output-preserved ceiling and the vertical line the measured retained wage share. Scenario, on the
assembled rate.}
\label{fig:boundary}
\end{figure}

\subsection{What the rate is a statement about}

Under full expensing the deduction precedes the income. The rate assembled here is therefore a
present-value wedge over the life of one marginal investment, while the replacement question asks
about tax collected in a year. The two are comparable only under three conditions: investment
constant, so that this year's deduction matches the deductions behind the income now taxed; rates
constant, so that the wedge on a project placed today is the wedge that applied to projects
already in place; and the capital stock stationary, so that both annual flows are steady. Under
those conditions the annual rate and the marginal wedge stand in the same ratio, and outside them
every comparison in this section fails. The transition runs against the paper rather than for it:
while investment grows, current deductions exceed the tax on income from a smaller earlier capital
stock, so revenue actually collected is below the assembled rate implies, and
Appendix~\ref{app:detail} sizes that with an illustrative path.

\subsection{Why the rate is low, and ownership matters more than profit shifting}
\label{sec:levers}

The binding constraint is who owns the claims. About \result{ThetaUntaxable} percent of US
corporate equity is held in forms the individual income tax barely reaches, chiefly retirement
accounts, foreign investors and nonprofits, leaving \result{ThetaTaxable} percent in taxable
accounts across a sourced range of \result{ThetaTaxableLow} to \result{ThetaTaxableHigh} percent
\citep{RosenthalAustin2016,RosenthalBurke2020,RosenthalMucciolo2024}. Of the gains that do accrue
there, \result{GainsUntaxedAtDeath} percent are held until death and never taxed, on 2014 data
\citep{CBO2014}, and what is realized is realized late, deferral cutting the present value by a
factor of \result{DeferralLow} to \result{DeferralHigh} on the later CRS estimates
\citep{Gravelle2022}. The two vintages are a decade apart and agree that roughly half of accrued
gains never bear tax, which is why we carry both rather than the newer alone. The debt side has the same
structure: \result{InterestSheltered} percent of corporate interest goes to holders who do not pay
tax on it and \result{BondsRestOfWorld} percent of corporate bonds are held abroad. Together these
are why the owner's layer adds only \result{ShareholderLayerPct} percent to the rate on rents.

Profit shifting is real. How much it matters is a different question from how uncertain it is,
and the two are easy to conflate.

\paragraph{Sensitivity is not importance, and the lever table measures sensitivity.}
Table~\ref{tab:levers} moves each component across its range with the others held at central
values. That reports how much the answer would move if we were wrong about a parameter. It does
not report how much that channel reduces the rate. A parameter pinned to a narrow range by good
evidence has a small swing and may still be the largest determinant of the level, which is
exactly the argument Appendix~\ref{app:detail} already makes about the taxable-holder share. The
same logic applies to every row, so the table is reported here as an uncertainty diagnostic and
no claim about relative importance is drawn from it.

\paragraph{The counterfactual comparison, by mechanism rather than by package.} To ask which
channel matters we set each to its no-friction value and measure what happens to the level.
Ownership is not one mechanism but two, the share of equity held in taxable accounts and the
deferral and step-up that erode what reaches those accounts, and bundling them would overstate
the comparison. They are kept apart.

\begin{table}[htbp]
\centering
\caption{Counterfactual contribution of each mechanism to the assembled marginal rate,
federal-only basis, Barkai rent reading. ``Alone'' sets that mechanism to its no-friction value
with the others at sourced centrals. The Shapley column averages each mechanism's marginal
contribution over all six orders of removal and sums exactly to the joint effect. Scenario.}
\label{tab:counterfactual}
\begin{tabular}{lrr}
\toprule
Mechanism set to its no-friction value & Alone & Shapley \\
\midrule
Taxable shareholder base complete & \result{ShThetaAlone} & \result{ShThetaShapley} \\
Deferral and step-up eliminated & \result{ShDeferAlone} & \result{ShDeferShapley} \\
Profit shifting eliminated & \result{ShShiftAlone} & \result{ShShiftShapley} \\
\midrule
All three together & & \result{ShJoint} \\
\bottomrule
\end{tabular}
\end{table}

Read alone, the three are closer than a bundled comparison suggests. Eliminating profit shifting
adds \result{ShShiftAlone} percentage points and eliminating deferral adds
\result{ShDeferAlone}, which is larger but not by much. The taxable shareholder base is the
outlier at \result{ShThetaAlone}.

The alone column understates the ownership mechanisms because they are complements: fixing
deferral makes the taxable base matter more and the reverse. The Shapley attribution, which
averages over every order of removal and sums exactly to the joint effect of
\result{ShJoint} points, puts the taxable base at \result{ShThetaShapley}, deferral at
\result{ShDeferShapley} and profit shifting at \result{ShShiftShapley}.

So the comparison the title draws is this one. The single largest mechanism, the taxable
shareholder base, is about \result{ShRatioTheta} times profit shifting. The two ownership
mechanisms together are about \result{ShRatioPair} times it, and that figure is a package and is
labelled as one. What does not hold is a claim that every ownership mechanism individually dwarfs
shifting: deferral does not.

Neither counterfactual is a policy. No reform abolishes profit shifting or places all equity in
taxable hands with no deferral, so the rows bound what each mechanism can contribute rather than
proposing anything.

Under the Karabarbounis and Neiman rent reading profit shifting contributes exactly nothing, and
that follows from the model rather than from the data: shifting operates only on the rent
component and that reading sets the rent share to zero. It is reported for completeness and
carries no evidential weight in the comparison.

\paragraph{The denominator caveat, which runs against us.} Both sources for the shifted share
measure a share of \emph{foreign} profits, while the assembly applies the parameter to the whole
rent base, treating a dollar of US rent as booked abroad with that probability. That overstates
shifting unless essentially all rent is earned through foreign affiliates. Correcting it would
raise the assembled rate and narrow the shortfall this paper reports, and it would also
\emph{shrink} the profit-shifting counterfactual above, which is the direction that strengthens
the comparison rather than the paper's headline. It is stated at L11 in
Section~\ref{sec:limitations} and bounded there.
\begin{table}[htbp]
\centering
\caption{The levers, ranked by sensitivity to parameter uncertainty rather than by importance.
Each is moved across its whole sourced or swept range with the others at central values, so a row
reports how much the answer would move if we were wrong about that parameter, not how much the
channel reduces the rate. For the latter see Table~\ref{tab:counterfactual}. Scenario.}
\label{tab:levers}
\input{tables/tab_levers.tex}
\end{table}


\paragraph{The separate checks are not cumulative, so here they are jointly.} This paper carries
two parameters that move the answer in the paper's own favour when they move: the retained wage
share $R$, whose specifications span \result{RSensLow} to \result{RSensHigh} and which lowers the
required rate as it rises, and the foreign-rent share $\phi$ of
Section~\ref{sec:limitations}, which raises the assembled rate as it falls. Checking them one at
a time and reporting both as though the checks accumulated would be a mistake, because the favourable ends compound.

\begin{table}[htbp]
\centering
\caption{The replacement condition over the retained wage share and the foreign-rent share
jointly, federal-only basis, narrower labor tax reading. Cells are the assembled rate less the
required rate in percentage points; negative is a shortfall. Scenario.}
\label{tab:joint}
\begin{tabular}{lrrrr}
\toprule
 & assembled rate & $R=0.568$ & $R=0.600$ & $R=0.637$ \\
\midrule
$\phi = 1.00$ & 7.75 & $-3.26$ & $-2.45$ & $-1.50$ \\
$\phi = 0.60$ & 8.30 & $-2.71$ & $-1.90$ & $-0.96$ \\
$\phi = 0.30$ & 8.71 & $-2.30$ & $-1.49$ & $-0.54$ \\
$\phi = 0.00$ & 9.13 & $-1.89$ & $-1.08$ & $-0.13$ \\
\bottomrule
\end{tabular}
\end{table}

The condition closes in \result{JointClosing} of the \result{JointCells} cells, so the direction
survives the joint move. The margin does not. At the corner most favourable to us, $\phi = 0$ and
$R = \result{RSensHigh}$, the shortfall is \result{JointNarrowest} percentage points against
\result{JointWidest} at the central case. That is close enough to nothing that we will not claim
the shortfall is established at that corner, only that it has not reversed anywhere on the grid.

Two things follow and both weaken the paper rather than strengthen it. The headline shortfall of
\result{JointWidest} points is a central-case figure and not a lower bound. And a reader who
holds different views about $R$ and $\phi$ than we do can reach a near-tie without leaving the
range either parameter is defensibly in. The result we are willing to defend is the sign across
this grid, not the size at any point in it.

\subsection{Checks, and what this module does not reach}

Two published figures check the assembly and one limitation runs against our own finding; both are set out in Appendix~\ref{app:detail}. In short, our shareholder layer corroborates the Congressional Research Service's text figure once the two are put on a common base, and the base construction treats foreign holders as bearing no tax at the owner level though dividends to them bear withholding tax. That effect is sized in Section~\ref{sec:fiscal} above: at a fifteen percent treaty rate it raises the assembled rate to \result{TauKWithholding} percent, and to \result{TauKJointWeakening} percent with the filers' own debt share as well, both below the \result{RequiredAllGovEasy} percent the condition requires. What stays open is the actual effective withholding rate, which varies by treaty and with the split between dividends, which bear it, and capital gains realized by foreign holders, which generally do not.
